\chapter{Objectivity}
\label{chapter:objectivity}

\section{Introduction}
\begin{itemize}
\item For two reference frames, $\cF$ and $\cF^*$, the base vectors
for $\cF$, $\vefa$, $\vefb$, and $\vefc$ transform through a rigid
body rotation $\tRbold$, such that 
\be
\vefa^* = \tRbold \vefa, \;\;\;
\vefb^* = \tRbold \vefb, \;\;\; \mbox{and} \;\;\;
\vefc^* = \tRbold \vefc.
\ee
%
\item A vector quantity $\va$ is objective if 
\be
 \va^* = \tRbold \va.
\ee
%
\item A tensor quantity $\tabold$ is objective if
\be
 \tabold^* = \tRbold \tabold \tRbold\transpose.
\ee
\end{itemize}

\begin{Hresult}
The Cauchy stress $\ve{\sigma}$ is objective.
\end{Hresult}

\begin{Hproof}
The traction vector $\vt$ is the inner product of the 
Cauchy stress tensor $\ve{\sigma}$ and the outward surface
normal vector $\vnhat$.
The same traction vector in a different reference frame is
$\vt^* = \ve{\sigma}^* \vnhat^*$.
Since vectors transform between reference frames by through 
a rotation tensor $\tRbold$, 
 $\vt^* = \tRbold \vt$
and
 $\vnhat^* = \tRbold \vnhat$.
Substituting these two vector transformations in the second
observer's frame of reference gives
\be
\tRbold \vt = \ve{\sigma}^* \tRbold \vnhat
\ee
but $\tRbold \vt = \tRbold \ve{\sigma} \vnhat$, so
$\tRbold \ve{\sigma} \vnhat = \ve{\sigma}^* \tRbold \vnhat$ for all
$\vnhat$, thus
\be
\ve{\sigma}^* = \tRbold \ve{\sigma} \tRbold\transpose,
\ee
which is how second-order tensors transform when they are objective.
Thus $\ve{\sigma}$ is objective.
\end{Hproof}


\section{Two-Point Tensors Transform as Vectors}
\begin{itemize}

\item The deformation gradient is a two-point tensor, 
with one coordinate in the reference configuration
and the other in the current configuration.  
The reference configuration
is inherently independent since it is tied to the observer.  
The current configuration makes the deformation gradient 
transform as a vector, not a tensor.  
For more information, see Holzapfel at 189.

\item The rotation is a two-point tensor, 
moving between the reference configuration and 
the current configuration.
The rotation thus transforms as a vector, not a tensor.  
For more information, see Holzapfel at 190.

\item The first Piola-Kirchhoff
stress tensor is a two-point tensor.  
Like the deformation gradient and
rotation, mentioned above, PK1 transforms 
as a vector, not a tensor.
See Holzapfel at 191 for more detail.

\end{itemize}

\section{Convected Rates of Change}

\section{Corotated Rates of Change}
An objective vector $\va$ 
\be
 \va^* = \tRbold \; \va,
\ee
and an objective tensor quantity $\tabold$ 
\be
 \tabold^* = \tRbold \; \tabold \; \tRbold\transpose,
\ee
generally do {\bf not} maintain their objectivity through
time differentiation, viz.,
\be
 \dot{\va}^* = \dot{\tRbold} \; \va + \tRbold \; \dot{\va},
\ee
and
\be
 \dot{\tabold}^* = 
 \dot{\tRbold} \; \tabold \; \tRbold\transpose + 
 \tRbold \; \dot{\tabold} \; \tRbold\transpose + 
 \tRbold \; \tabold \; \dot{\tRbold}\transpose.
\ee

But, material vector fields can be endowed with objectivity.
For vector $\va$, the 
corotational rate 
% $\mathring{\va}$ % former notation
$\overset{\triangledown}{\va}$ is defined as
\be
 \overset{\triangledown}{\va}
 \defe \dot{\va} - \twbold \; \va,
\ee
where $\twbold$ is the rotation rate.

For a second-order tensor $\tabold$, the 
corotational rate $\overset{\triangledown}{\tabold}$ is defined as
\be
 \overset{\triangledown}{\tabold} \defe \dot{\tabold} - \twbold \; \tabold
 + \tabold \; \twbold,
\ee
which is referred to as the {\bf Jaumann-Zaremba rate}.

